% =====================================================================
% Madhur N Patel — résumé
%
% SINGLE SOURCE OF TRUTH for public/Madhur_N_Patel_Resume.pdf.
% Build with:  npm run resume     (see resume/build.sh)
%
% Content is mirrored from the site's data files:
%   src/data/profile.ts       header, summary
%   src/data/projects.ts      the three flagship projects + repo URLs
%   src/data/experience.ts    experience bullets (verbatim)
%   src/data/skills.ts        skills
%   src/data/certifications.ts
%
% Where this file and the site disagree, THE SITE WINS — it is the version
% that was deliberately cleaned. Update the data file, then this, then rebuild.
%
% Never reintroduce: "3+ years", "95%", "100+ concurrent users",
% "zero-downtime", "Kubernetes", "Freelance", "Expert". These were removed
% from the site for cause (unverifiable or inaccurate) and must not come back
% through the résumé.
%
% Keep this to ONE PAGE. If it overflows, cut the "Additional Projects" line
% and then the older certifications — never shrink below 10pt or margins
% below 0.5in.
% =====================================================================
\documentclass[10pt,a4paper]{article}

\usepackage[T1]{fontenc}
\usepackage[utf8]{inputenc}
\usepackage[margin=0.5in]{geometry}
\usepackage{titlesec}
\usepackage{enumitem}
\usepackage[dvipsnames]{xcolor}
\usepackage[hidelinks]{hyperref}

\hypersetup{
  colorlinks=true,
  urlcolor=[HTML]{1A5FB4},
  linkcolor=[HTML]{1A5FB4},
  pdftitle={Madhur N Patel — Infrastructure \& Embedded Systems Engineer},
  pdfauthor={Madhur N Patel},
  pdfsubject={Resume},
  pdfkeywords={infrastructure, embedded systems, Linux, Docker, Raspberry Pi, ESP32}
}

% --- Layout -----------------------------------------------------------
\pagestyle{empty}
\setlength{\parindent}{0pt}
\setlength{\tabcolsep}{0pt}

\titleformat{\section}
  {\large\bfseries\scshape}
  {}{0pt}{}[\vspace{-0.7em}\rule{\linewidth}{0.6pt}]
\titlespacing*{\section}{0pt}{0.55em}{0.35em}

\newcommand{\entry}[2]{%
  \textbf{#1}\hfill{\small #2}\par\vspace{0.15em}}

\newlist{tight}{itemize}{1}
\setlist[tight]{leftmargin=1.1em, itemsep=0.05em, topsep=0.1em, parsep=0pt, label=\textbullet}

\begin{document}

% --- Header -----------------------------------------------------------
\begin{center}
  {\huge\bfseries Madhur N Patel}\\[0.35em]
  {\large Infrastructure \& Embedded Systems Engineer}\\[0.5em]
  {\small
    Ahmedabad, Gujarat, India \quad\textbar\quad
    \href{mailto:pmadhurn@gmail.com}{pmadhurn@gmail.com} \quad\textbar\quad
    +91 9016273812\\
    \href{https://madhur.dev}{madhur.dev} \quad\textbar\quad
    \href{https://github.com/pmadhurn}{github.com/pmadhurn} \quad\textbar\quad
    \href{https://linkedin.com/in/madhur-n}{linkedin.com/in/madhur-n}
  }
\end{center}

\vspace{-0.4em}

% --- Summary ----------------------------------------------------------
\section{Summary}
Builds and operates self-managed Linux infrastructure, containerized services, and
sensor-fusion systems on Raspberry Pi and ESP32. Looking for infrastructure, embedded,
or platform engineering work — ideally somewhere that ships physical products.

% --- Projects (above education: these are the strongest evidence) ------
\section{Projects}

% Repo URL removed 2026-08-15 — github.com/pmadhurn/gps-guided-tracking-system 404s.
\entry{GPS-Guided Dual-Device Tracking System}{Python · OpenCV · Raspberry Pi}
\begin{tight}
  \item Two Raspberry Pi--controlled gimbals point toward each other by GPS bearing, then
        hand off to computer vision for fine LED tracking; automatic mode switching between
        coarse GPS and camera tracking.
  \item Real-time 9-axis IMU650 sensor fusion (GPS, IMU, magnetometer), Haversine distance
        and elevation-angle calculation, SP2050 gimbal control under PID regulation.
  \item Multi-threaded serial and socket communication between the two devices over WiFi.
  \item Python, OpenCV, PySerial, Raspberry Pi 4/5.
\end{tight}

\vspace{0.25em}
% Repo URL removed 2026-08-15 — pmadhurn/OCI is an unrelated capacity script, not this project.
\entry{Self-Hosted Infrastructure — OCI ARM64 + Cloudflare Tunnel}{Docker · Cloudflare Tunnel}
\begin{tight}
  \item Self-managed production environment running six containerized services on a
        4-core / 24\,GB ARM64 Oracle Cloud instance (Open WebUI, n8n, Portainer,
        code-server, Ollama, Supabase).
  \item Zero-trust ingress through Cloudflare Tunnel with no public inbound ports;
        subdomain routing via tunnel ingress rules with an ordered catch-all fallback.
  \item Local LLM inference on ARM64 via Ollama; self-hosted Supabase with credential
        isolation; service monitoring, backup, and TLS certificate management.
  \item Docker, Docker Compose, Nginx, Linux (Ubuntu), PostgreSQL.
\end{tight}

\vspace{0.25em}
\entry{SpeakInsights — Meeting Intelligence Platform}{\href{https://github.com/pmadhurn/SpeakinsightsV6}{github.com/pmadhurn/SpeakinsightsV6}}
\begin{tight}
  \item Docker-deployable meeting platform: LiveKit WebRTC conferencing for up to 20
        participants, WhisperX transcription with speaker attribution, Ollama summarization
        and task extraction, RAG chat over meeting context with pgvector.
  \item Runs fully on-premises — no third-party API calls, so meeting audio never leaves
        the deployment.
  \item Async FastAPI backend with a Redis-backed job queue; batched GPU inference for
        long-form audio; rolling container restarts via Docker Compose.
  \item React 18, TypeScript, FastAPI, Python 3.11, PostgreSQL + pgvector, Redis, Docker.
\end{tight}

% --- Experience -------------------------------------------------------
\section{Experience}

\entry{Self-Directed Infrastructure Engineering}{2022 -- Present}
\begin{tight}
  \item Implemented robust CI/CD pipelines using GitHub Actions, automating the build,
        test, and deployment of applications.
  \item Engineered a secure and scalable network infrastructure using Cloudflare Tunnel,
        enabling secure access to internal applications and services without exposing them
        to the public internet.
  \item Managed Vercel deployments with a non-proxied Cloudflare setup, optimizing for
        performance and security by leveraging Cloudflare's DNS and security features while
        directly serving content from Vercel.
  \item Containerized services with Docker and Docker Compose, managing service
        dependencies, health checks, and persistent volumes across a six-service deployment.
\end{tight}

% --- Education --------------------------------------------------------
\section{Education \& Certifications}
\entry{B.Tech, Computer Science and Engineering — Indus University, Ahmedabad}{2022 -- 2026}
GPA 8.5/10\\[0.25em]
AWS Academy Graduate — Machine Learning Foundations (Aug 2025) \textbullet{}
AWS Academy Graduate — Cloud Foundations (Jul 2025) \textbullet{}
MERN Stack Developer, Coding Cloud (Apr 2025)

% --- Skills -----------------------------------------------------------
% Flowing lines, not a tabular: a tabular is one unbreakable box, and if it
% doesn't fit the remaining space the whole section jumps to page 2.
\section{Skills}
\setlength{\parskip}{0.1em}%
\textbf{Strongest:} Python, JavaScript, Node.js, React.js, Git/GitHub

\textbf{Infrastructure:} Docker, Docker Compose, Cloudflare Tunnel, Nginx, Oracle Cloud, Vercel, CI/CD, AWS, GCP

\textbf{Embedded/IoT:} Raspberry Pi, ESP32, Arduino, Serial / I2C \quad\textbullet\quad \textbf{Systems:} Linux (Ubuntu), macOS, Windows

\textbf{Languages:} Python, JavaScript, TypeScript, SQL, C, C++, Java, Bash

\textbf{Backend:} FastAPI, Node.js, Express.js, REST APIs \quad\textbullet\quad \textbf{Databases:} PostgreSQL, MongoDB, MySQL, SQLite, Redis

\textbf{AI/ML:} OpenCV, Ollama (local inference), Pandas, NumPy, TensorFlow, scikit-learn

\end{document}
